%%%%%%%%%%%%%%%%%%%%%%%%%%%%%%%%%%%%%%%%%%%%%%%%%%%%%%%%%%%%%%%%%%
% Abstract em Inglês
%%%%%%%%%%%%%%%%%%%%%%%%%%%%%%%%%%%%%%%%%%%%%%%%%%%%%%%%%%%%%%%%%%

\chapter*{Abstract}
\addcontentsline{toc}{chapter}{Abstract}

% Embedded hyperspectral processing represents a significant challenge due to
% energy and latency constraints imposed by mobile and autonomous systems. This
% dissertation presents a systematic study on computational optimization
% strategies for heterogeneous systems (CPU+GPU+FPGA) applied to real-time
% hyperspectral processing.

% The research was structured in two distinct phases: Phase 1 consists of
% systematic state-of-the-art analysis, cataloging 25 scientific works and
% identifying proven optimization techniques; Phase 2, planned for
% post-qualification, will involve experimental development and practical
% validation of proposed solutions.

% The systematic analysis revealed significant gaps in current literature: 96\% of
% works focus on isolated techniques, 88\% report partial metrics, and 80\% use
% synthetic validation. As the main contribution of Phase 1, a conceptual
% architectural framework was developed that integrates compressive sensing
% techniques, EMCR (Extreme Learning Machine with Correlation Removal) spectral
% selection, and hardware/software codesign.

% The identified optimal configuration consists of: (1) Zynq UltraScale+ FPGA for
% correction and EMCR selection, providing 80\% data reduction with <5ms latency
% and <5W consumption; (2) Jetson Orin GPU for compressive sensing/PCA processing,
% offering 50-70\% additional reduction with <10ms and <10W; (3) ARM Cortex-A78
% CPU for final classification using lightweight CNNs or SVM, maintaining <5ms and
% <5W.

% Theoretical results indicate potential to achieve throughput >100 fps, total
% latency <50ms, consumption <15W, and accuracy >95\%. The estimated speedup is
% 18-21x compared to CPU baseline, with energy efficiency superior to 20x. These
% metrics establish a solid foundation for Phase 2 experimental validation,
% contributing to the advancement of embedded hyperspectral processing in critical
% applications such as precision agriculture, environmental monitoring, and
% autonomous UAV systems.

\textbf{Keywords:} 
% Hyperspectral Processing, Heterogeneous Systems, Energy Optimization, Compressive Sensing, Hardware/Software Codesign, Embedded Systems.
